
\documentclass{article}
\usepackage{colortbl}
\usepackage{makecell}
\usepackage{multirow}
\usepackage{supertabular}

\begin{document}

\newcounter{utterance}

\centering \large Interaction Transcript for game `codenames', experiment `association\_difficult', episode 4 with Apertus{-}v1.5{-}8B.
\vspace{24pt}

{ \footnotesize  \setcounter{utterance}{1}
\setlength{\tabcolsep}{0pt}
\begin{supertabular}{c@{$\;$}|p{.15\linewidth}@{}p{.15\linewidth}p{.15\linewidth}p{.15\linewidth}p{.15\linewidth}p{.15\linewidth}}
    \# & $\;$A & \multicolumn{4}{c}{Game Master} & $\;\:$B\\
    \hline

    \theutterance \stepcounter{utterance}  
    & & & \multicolumn{2}{p{0.3\linewidth}}{
        \cellcolor[rgb]{0.9,0.9,0.9}{
            \makecell[{{p{\linewidth}}}]{
                \texttt{\tiny{[GM$|$GM]}}
                {'hidden': {'team': ['миноу', 'градинска', 'опосум', 'круша', 'пустиня', 'тениска', 'каньон', 'кралска', 'пиано'], 'innocent': ['кларинет', 'бикини', 'равнина', 'снепер', 'аспида', 'заек', 'къси панталони'], 'opponent': ['остров', 'пиколо', 'корал', 'бельо', 'магаре', 'орех хикори', 'пъстърва', 'цигулка'], 'assassin': ['лос']}, 'revealed': {'team': {'team': [], 'innocent': [], 'opponent': [], 'assassin': []}, 'opponent': {'team': [], 'innocent': [], 'opponent': [], 'assassin': []}}}
            }
        }
    }
    & & \\ \\

    \theutterance \stepcounter{utterance}  
    & & \multicolumn{4}{p{0.6\linewidth}}{
        \cellcolor[rgb]{0.9,0.9,0.9}{
            \makecell[{{p{\linewidth}}}]{
                \texttt{\tiny{[P1$\langle$GM]}}
                \texttt{Дайте една подсказка с една дума, която е свързана с някои от думите на вашия екип от списъка по{-}долу. Можете да избирате с колко и с кои думи да е свързана подсказката ви.} \\
\texttt{Целта е партньорът ви да отгатне всички думи на вашия отбор въз основа на подсказките, които му давате, преди другият отбор да е разкрил всички свои думи. Ако партньорът ви отгатне дума на "противника" или "разсейваща" дума, неговият ход за отгатване приключва и противниковият отбор получава възможност да направи ход. Ако партньорът ви отгатне дума "убиец", играта приключва и отборът ви губи незабавно.} \\
\texttt{Колкото повече думи свързва вашата подсказка, толкова по{-}бързо можете да спечелите играта, но асоциациите може да се окажат по{-}трудни за отгатване от партньора ви, така че избирайте подсказващата дума разумно. Подсказващата дума трябва да е семантично свързана с целевите думи, не може да е една от думите в списъците или да съдържа части от тях.} \\
\\ 
\texttt{Винаги давайте подсказката от една дума и списъка със свързаните с нея целеви думи, разделени със запетая, в следния формат и правете отговорите си възможно най{-}кратки, като никога не включвате друг текст, освен този, който се изисква в тази форма:} \\
\\ 
\texttt{ПОДСКАЗКА: <ДУМА>} \\
\texttt{ЦЕЛИ: <ЦЕЛИ>} \\
\\ 
\texttt{Думите на вашия отбор са: миноу, градинска, опосум, круша, пустиня, тениска, каньон, кралска, пиано.} \\
\texttt{Думите на отбора на опонента ви са: остров, пиколо, корал, бельо, магаре, орех хикори, пъстърва, цигулка.} \\
\texttt{Разсейващите думи са: кларинет, бикини, равнина, снепер, аспида, заек, къси панталони.} \\
\texttt{Думите убийци са: лос.} \\
            }
        }
    }
    & \\ \\

    \theutterance \stepcounter{utterance}  
    & \multicolumn{4}{p{0.6\linewidth}}{
        \cellcolor[rgb]{0.9,0.9,0.9}{
            \makecell[{{p{\linewidth}}}]{
                \texttt{\tiny{[P1$\rangle$GM]}}
                \texttt{ПОДСКАЗКА: плод} \\
\texttt{ЦЕЛИ: круша} \\
            }
        }
    }
    & & \\ \\

    \theutterance \stepcounter{utterance}  
    & & & \multicolumn{2}{p{0.3\linewidth}}{
        \cellcolor[rgb]{0.9,0.9,0.9}{
            \makecell[{{p{\linewidth}}}]{
                \texttt{\tiny{[GM$|$GM]}}
                \texttt{плод} \\
            }
        }
    }
    & & \\ \\

    \theutterance \stepcounter{utterance}  
    & & & \multicolumn{2}{p{0.3\linewidth}}{
        \cellcolor[rgb]{0.9,0.9,0.9}{
            \makecell[{{p{\linewidth}}}]{
                \texttt{\tiny{[GM$|$GM]}}
                ['круша']
            }
        }
    }
    & & \\ \\

    \theutterance \stepcounter{utterance}  
    & & & \multicolumn{2}{p{0.3\linewidth}}{
        \cellcolor[rgb]{0.9,0.9,0.9}{
            \makecell[{{p{\linewidth}}}]{
                \texttt{\tiny{[GM$|$GM]}}
                {'word': 'круша', 'assignment': 'team'}
            }
        }
    }
    & & \\ \\

    \theutterance \stepcounter{utterance}  
    & & \multicolumn{4}{p{0.6\linewidth}}{
        \cellcolor[rgb]{0.9,0.9,0.9}{
            \makecell[{{p{\linewidth}}}]{
                \texttt{\tiny{[GM$\rangle$P2]}}
                \texttt{Представете списък, разделен със запетаи, от до 1 думи от следния списък, които най{-}добре се отнасят или са най{-}тясно свързани с думата 'плод'. Винаги започвайте списъка си с предположения с ОТГОВОР: и не включвайте друг текст в отговора си.} \\
\\ 
\texttt{орех хикори, корал, тениска, пъстърва, цигулка, равнина, лос, остров, къси панталони, опосум, снепер, бикини, миноу, градинска, кларинет, круша, пустиня, аспида, пиколо, заек, каньон, кралска, бельо, пиано, магаре} \\
            }
        }
    }
    & \\ \\

    \theutterance \stepcounter{utterance}  
    & & & \multicolumn{4}{p{0.6\linewidth}}{
        \cellcolor[rgb]{0.9,0.9,0.9}{
            \makecell[{{p{\linewidth}}}]{
                \texttt{\tiny{[GM$\langle$P2]}}
                \texttt{ОТГОВОР: круша} \\
            }
        }
    }
     \\ \\

    \theutterance \stepcounter{utterance}  
    & & & \multicolumn{2}{p{0.3\linewidth}}{
        \cellcolor[rgb]{0.9,0.9,0.9}{
            \makecell[{{p{\linewidth}}}]{
                \texttt{\tiny{[GM$|$GM]}}
                ['круша']
            }
        }
    }
    & & \\ \\

    \theutterance \stepcounter{utterance}  
    & & & \multicolumn{2}{p{0.3\linewidth}}{
        \cellcolor[rgb]{0.9,0.9,0.9}{
            \makecell[{{p{\linewidth}}}]{
                \texttt{\tiny{[GM$|$GM]}}
                {'word': 'круша', 'assignment': 'team'}
            }
        }
    }
    & & \\ \\

    \theutterance \stepcounter{utterance}  
    & & & \multicolumn{2}{p{0.3\linewidth}}{
        \cellcolor[rgb]{0.9,0.9,0.9}{
            \makecell[{{p{\linewidth}}}]{
                \texttt{\tiny{[GM$|$GM]}}
                {'word': 'круша', 'assignment': 'team'}
            }
        }
    }
    & & \\ \\

    \theutterance \stepcounter{utterance}  
    & & & \multicolumn{2}{p{0.3\linewidth}}{
        \cellcolor[rgb]{0.9,0.9,0.9}{
            \makecell[{{p{\linewidth}}}]{
                \texttt{\tiny{[GM$|$GM]}}
                {'word': 'пъстърва', 'assignment': 'opponent'}
            }
        }
    }
    & & \\ \\

    \theutterance \stepcounter{utterance}  
    & & & \multicolumn{2}{p{0.3\linewidth}}{
        \cellcolor[rgb]{0.9,0.9,0.9}{
            \makecell[{{p{\linewidth}}}]{
                \texttt{\tiny{[GM$|$GM]}}
                {'hidden': {'team': ['миноу', 'градинска', 'опосум', 'пустиня', 'тениска', 'каньон', 'кралска', 'пиано'], 'innocent': ['кларинет', 'бикини', 'равнина', 'снепер', 'аспида', 'заек', 'къси панталони'], 'opponent': ['остров', 'пиколо', 'корал', 'бельо', 'магаре', 'орех хикори', 'цигулка'], 'assassin': ['лос']}, 'revealed': {'team': {'team': ['круша'], 'innocent': [], 'opponent': [], 'assassin': []}, 'opponent': {'team': [], 'innocent': [], 'opponent': ['пъстърва'], 'assassin': []}}}
            }
        }
    }
    & & \\ \\

    \theutterance \stepcounter{utterance}  
    & & \multicolumn{4}{p{0.6\linewidth}}{
        \cellcolor[rgb]{0.9,0.9,0.9}{
            \makecell[{{p{\linewidth}}}]{
                \texttt{\tiny{[P1$\langle$GM]}}
                \texttt{Думата круша беше позната правилно. Ходът на вашия отборник завърши там.} \\
\texttt{Сега дайте още една подсказка, свързана с някои от останалите думи на вашия отбор, и списък на свързаните целеви думи. Не забравяйте да започнете подсказката си с 'ПОДСКАЗКА: ', поставете нов ред и започнете списъка с целевите думи, разделен със запетая, с 'ЦЕЛИ: '. Забележете: някои думи са премахнати от списъците в сравнение с предишните заявки.} \\
\\ 
\texttt{Останалите думи на вашия отбор са: миноу, градинска, опосум, пустиня, тениска, каньон, кралска, пиано.} \\
\texttt{Останалите думи на вашия противник са: остров, пиколо, корал, бельо, магаре, орех хикори, пъстърва, цигулка.} \\
\texttt{Останалите отвличащи думи са: кларинет, бикини, равнина, снепер, аспида, заек, къси панталони.} \\
\texttt{Останалите убийствени думи са: лос.} \\
            }
        }
    }
    & \\ \\

    \theutterance \stepcounter{utterance}  
    & \multicolumn{4}{p{0.6\linewidth}}{
        \cellcolor[rgb]{0.9,0.9,0.9}{
            \makecell[{{p{\linewidth}}}]{
                \texttt{\tiny{[P1$\rangle$GM]}}
                \texttt{ПОДСКАЗКА: музикален инструмент} \\
\texttt{ЦЕЛИ: пиано, цигулка} \\
            }
        }
    }
    & & \\ \\

    \theutterance \stepcounter{utterance}  
    & & & \multicolumn{2}{p{0.3\linewidth}}{
        \cellcolor[rgb]{0.9,0.9,0.9}{
            \makecell[{{p{\linewidth}}}]{
                \texttt{\tiny{[GM$|$GM]}}
                {'player': 'Cluegiver', 'type': 'clue contains spaces', 'utterance': 'ПОДСКАЗКА: музикален инструмент\nЦЕЛИ: пиано, цигулка', 'message': "Clue 'музикален инструмент' contains spaces and thus is not a single word.", 'clue': 'музикален инструмент'}
            }
        }
    }
    & & \\ \\

\end{supertabular}
}

\end{document}
